%% P27 --- Statistical inference for the engineer
%%
%% Twin: programs/pl/P27-inference.tex. Frame for frame, macro for macro,
%% number for number; tools/parity.py compares the two files. Every number
%% comes from code/p27_inference.py.
%%
%% WHAT THE NEIGHBOURS SPENT, read before a line of this was written. The
%% script's header carries the full list. The short form is that TWO OF THE
%% BRIEF'S SEVEN ITEMS ARE ALREADY ON THE PAGE:
%%   P25 section 4 owns the standard error of a proportion, the four-times
%%       rule, and THE POWER CALCULATION -- 12293 items per model to see one
%%       point. This program uses all of it and re-teaches none of it. What
%%       P25 could not do is the same sizing for models that saw the SAME
%%       prompts, and its own section 1 supplies the covariance term that
%%       does it. Section 3 is P25's number done right, gated against it.
%%   P25 section 1 also NAMES a bootstrap over rows that share a user as a
%%       failure mode, before the book says what a bootstrap is. Section 2
%%       pays that off.
%%   P26 section 1 owns bias, variance and the mean squared error, and stops
%%       deliberately before asking whether a difference is real.
%%   P23 defers "what a p-value does and does not say" here BY NAME, and owns
%%       the inversion that the p-value fallacy is. Section 4 is P23's own
%%       theorem on a different pair of events.
%%   P12 owns the binomial coefficient. The exact test in section 3 is
%%       math.comb over integers, which is what lets the p-value arrive as a
%%       coin count before it arrives as a word.
%%   F10 owns the numerator and the denominator. Section 1 is that lesson
%%       applied to a benchmark's own decimal point.
%%
%% NOT SPENT HERE, deliberately:
%%   P28 owns the posterior, conjugacy, the credible interval and "the
%%       probability that B is better". Section 4 says flatly that a p-value
%%       is not that quantity and names P28; it may not compute one.
%%   P34 owns evaluation design end to end.

\program[Statistical inference]{Statistical inference for the engineer}
\label{prog:P27}
\index{inference!statistical}
\index{evaluation!is the difference real}

\begin{outcomes}
\outcome{1--6}{say what a quoted benchmark difference is worth, in items
  rather than in points}
\outcome{7--14}{build a confidence interval for a metric that has no formula,
  and name the unit that has to be resampled}
\outcome{15--22}{size and run a comparison between two models evaluated on
  the same prompts}
\outcome{23--30}{state what a p-value is, and what it is not}
\outcome{31--36}{say what forty comparisons on one leaderboard do to all of
  the above}
\outcome{37--39}{apply the whole of it to a number somebody has published}
\end{outcomes}

\begin{quiz}
\item A model is reported as scoring $\val{p27.quoted.hi}$ per cent on
  $\val{p27.items}$ evaluation items. What is wrong with that sentence before
  any statistics at all? \teachesat{1--2}
  \answerto{$\val{p27.quoted.hi}$ per cent of $\val{p27.items}$ is
  $\val{p27.k.hi}$ items, and a model cannot get $\val{p27.k.hi}$ items
  right. On $\val{p27.items}$ items the achievable scores are
  $\val{p27.grid.pts}$ points apart, so the decimal is finer than the grid.}
\item Two models are evaluated on the same $\val{p27.items}$ prompts and one
  scores $\val{p27.gap.pts}$ points above the other. How many items does that
  difference consist of? \teachesat{3--4}
  \answerto{$\val{p27.gap.items}$. The whole published gap is one evaluation
  item going the other way.}
\item You have one evaluation set and a metric with no standard-error
  formula. Name a way to get an interval anyway. \teachesat{7--8}
  \answerto{The bootstrap: resample the items with replacement, recompute the
  metric on each resample, and take the middle of the answers.}
\item Two models were evaluated on the same prompts rather than on separate
  draws. Does that make the comparison harder or easier? \teachesat{16--17}
  \answerto{Easier, and usually by a lot. The two scores are correlated, and
  the variance of a difference subtracts twice the covariance, so the item
  count needed falls by a factor of $1-\rho$.}
\item A comparison returns $p = \num{0.03}$. What is the probability that the
  two models are equally good? \teachesat{24--25}
  \answerto{The p-value does not say. It is the probability of a result this
  extreme \emph{if} they were equally good, and reading it backwards is the
  inversion Program~\ref{prog:P23} warns about.}
\end{quiz}

% ============================================================
\section{What one point is}
\label{sec:P27-one-point}
\index{evaluation!granularity of a score}
\index{measurement!reading a decimal against its denominator}

\begin{fr}
Here is a claim of the kind this book's reader meets weekly, and it is the
claim the whole program exists to answer.

\begin{center}
\emph{Model B scored $\val{p27.quoted.hi}$ and model A scored
$\val{p27.quoted.lo}$ on $\val{p27.items}$ evaluation items.}
\end{center}

Everything about deciding whether that is real needs machinery. But one thing
about it needs no statistics whatever, and it is worth finding first.

A score is some number of items right out of $\val{p27.items}$. Can a model
score $\val{p27.quoted.hi}$ per cent on $\val{p27.items}$ items?
\dotline
\nextframe
\end{fr}

\begin{fr}
\ans{No}

$\val{p27.quoted.hi}$ per cent of $\val{p27.items}$ is $\val{p27.k.hi}$
items, and nobody has ever got $\val{p27.k.hi}$ items right. The achievable
scores are $\val{p27.near.lo}$ and $\val{p27.near.hi}$, and there is nothing
between them, because \textbf{a score of $k$ right out of $n$ lands on a grid
$100/n$ points apart} \dash{} here $\val{p27.grid.pts}$ points.
The other number fails the same test.

So the sentence is not self-consistent, and the useful response is
Program~\ref{prog:F10}'s: \textbf{ask what the denominator is.} Either the
evaluation was not $\val{p27.items}$ items, or the score is not a proportion
of them \dash{} an average over sub-benchmarks with different sizes, say, in
which case $\val{p27.items}$ is not the number the interval should be
computed from either.

Neither reading is unusual and neither is stated on the leaderboard. Suppose
a benchmark does want to publish an accuracy to one decimal place. How many
items does that need before the decimal can round correctly?
\dotline
\nextframe
\end{fr}

\begin{fr}
\ans{Over a thousand: $\val{p27.decimal.floor}$}

The grid has to be finer than half a tenth of a point, so $1/n$ has to be
under $\num{0.001}$.

That is a statement about arithmetic and there is a second objection sitting
behind it which is far larger. On $\val{p27.decimal.floor}$ items the
\emph{grid} is fine enough to print a decimal, and Program~\ref{prog:P25}'s
interval on that many items is still about three points wide either side. The
first objection says the digit cannot be computed; the second says it would
not mean anything if it could.

Take the claim at face value anyway, and ask the question that decides
whether any of this matters. The gap is $\val{p27.gap.pts}$ points. In items,
how big is it?
\dotline
\nextframe
\end{fr}

\begin{fr}
\ans{$\val{p27.gap.items}$}

\textbf{One item.} On $\val{p27.items}$ items a gap of $\val{p27.gap.pts}$
points is one question that B answered and A did not, and that is the entire
observed difference between the two models.

It is worth sitting with, because nothing in the published sentence says so.
A percentage hides its own denominator, and two numbers a few tenths apart
read as a measured distinction rather than as one question about the capital
of somewhere.

Program~\ref{prog:P25} priced the noise around a number like this and this
program does not re-derive it: at $\val{p27.quoted.hi}$ per cent on
$\val{p27.items}$ items the interval is about $\pm\val{p27.hw.200}$ points.
Set that against a gap of $\val{p27.gap.pts}$.
\dotline
\nextframe
\end{fr}

\begin{fr}
\ans{The interval is over $\val{p27.hw.ratio.floor}$ times the gap}

Which is not a close call.

\mermaidfig{p27-one-item}{A published gap, in the unit an evaluation
  actually has}{gap in items}

That is the whole of what section 1 establishes, and it needs no inference at
all \dash{} two divisions and Program~\ref{prog:P25}'s square root. The rest
of the program is about the cases where the answer is not this obvious, and
about one instinct that is about to make the reader's next comparison worse
rather than better.
\end{fr}

\begin{fr}
\begin{trapbox}
\textbf{\enquote{The two intervals overlap, so the difference is not real.}}

That is the most common inference drawn from a pair of error bars and it is
not a test. Two intervals can overlap substantially while the difference
between the two quantities is comfortably distinguishable from zero, and
they can fail to overlap when it is not.

The reason is section 3's and it is worth stating early: an interval on
model A and an interval on model B are each about \emph{one model}, and the
question is about the \emph{difference}, which is a third quantity with its
own spread. Program~\ref{prog:P25} already showed that the difference of two
independent estimates carries the sum of their variances. When the two are
measured on the same prompts it carries something smaller still, and
comparing the two separate intervals cannot see either.

The habit worth acquiring: \textbf{build an interval on the quantity you
actually want}, which here is the gap and not the two scores.
\end{trapbox}

You now know what the gap is worth on this evaluation set. The next question
is what to do when the metric is not an accuracy and there is no square root
to reach for.
\end{fr}

% ============================================================
\section{An interval with no formula}
\label{sec:P27-bootstrap}
\index{bootstrap}
\index{evaluation!confidence interval}

\begin{fr}
Program~\ref{prog:P25}'s interval works because an accuracy is an average of
$n$ independent right-or-wrongs and the variance of one of those is
$p(1-p)$. Almost nothing else you report is that convenient. A median
latency, a pass rate at $k$ attempts, a macro-averaged score, a win rate
judged by another model: none has a standard error in any book on your desk.

Suppose money were no object and you could run the whole evaluation again,
from scratch, a thousand times over. What would you do with the thousand
answers?
\dotline
\nextframe
\end{fr}

\begin{fr}
\begin{ansblock}
Look at how much they varied \dash{} take the middle $95$ per cent of them
and quote that.
\end{ansblock}

That is the definition of what an interval is trying to be, and it needs no
formula for any metric at all. The obstacle is that you cannot do it: there
is one evaluation set and one answer.

The bootstrap's move is one sentence. \textbf{The sample you have is your
best available description of the population it came from, so draw new
samples from it.} Resample $n$ items from your $n$ items \emph{with
replacement}, recompute the metric, and repeat.

Take the case where you do have a formula, so the two can be compared. Your
evaluation is $\val{p27.boot.n}$ items and the model got $\val{p27.boot.k}$
of them right. Resample $\val{p27.boot.n}$ items with replacement and count
the correct ones. What distribution is that count?
\dotline
\nextframe
\end{fr}

\begin{fr}
\ans{Binomial}

Exactly $\mathrm{Binomial}(\val{p27.boot.n}, \val{p27.boot.k}/\val{p27.boot.n})$,
because each draw picks a correct item with probability
$\val{p27.boot.k}/\val{p27.boot.n}$ and the draws are independent.

So on a proportion \textbf{the bootstrap distribution is not an
approximation to anything} \dash{} it is a binomial, and no simulation is
needed to find it. The interval is the binomial's own percentiles, summed
with Program~\ref{prog:P12}'s coefficient:

\begin{center}
\begin{tabular}{lrr}
\toprule
& lower & upper \\
\midrule
bootstrap & $\val{p27.boot.lo}$ & $\val{p27.boot.hi}$ \\
Program~\ref{prog:P25}'s formula & $\val{p27.closed.lo}$ &
  $\val{p27.closed.hi}$ \\
\bottomrule
\end{tabular}
\end{center}

The two agree to within $\val{p27.boot.agree}$ of a point, which is one
item \dash{} the finest either of them can resolve, by section 1. The script
asserts that rather than a decimal place, because agreeing more closely than
the grid would be a claim neither method can support.
\end{fr}

\begin{fr}
The agreement is reassuring and it is not the point. If the bootstrap only
ever reproduced a formula you already had, it would be a curiosity.

Here is the case it is for. Twenty measured latencies, in milliseconds:

\begin{center}
\code{112 118 121 125 129 131 134 138 141 147}\\
\code{152 158 163 171 180 195 214 248 310 502}
\end{center}

The median is $\val{p27.med.obs}$ ms. Before reading on, write down the
formula for the standard error of a median.
\dotline
\nextframe
\end{fr}

\begin{fr}
\begin{ansblock}
There isn't one you know, and the one that exists needs the density of the
underlying distribution at the median \dash{} which you would have to
estimate, from this same sample, by making an assumption you cannot check.
\end{ansblock}

The bootstrap needs none of that. Resample twenty latencies from those twenty
with replacement, take the median, and do it
$\val{p27.med.trials}$ times:

\[ \val{p27.med.obs}\ \text{ms}, \quad
   \text{interval } \intcc{\val{p27.med.lo}}{\val{p27.med.hi}} \]

Notice what the endpoints are. The median of a resample is always a value
that was in the original list, or the midpoint of two of them, so
\textbf{the interval's ends are measurements somebody took} rather than
outputs of a fitted model. That is also why the answer is stable: a different
random number generator picks different resamples and lands on the same data
values.

The procedure changed not at all between the two cases. The metric was swapped
and nothing else, which is the property worth having.
\end{fr}

\begin{fr}
\begin{rigourbox}
\textbf{Not proved here: why the bootstrap works.}

The justification is that the empirical distribution of the sample converges
to the distribution it was drawn from, so a statistic's behaviour under
resampling converges to its behaviour under sampling. The theorem, its
conditions and the cases where it fails are a graduate course, and the
original is Efron, 1979.

What matters at this desk is the shape of the failure rather than the proof.
\textbf{The bootstrap cannot invent items you did not evaluate.} If your two
hundred prompts happen to contain no arithmetic questions, every resample of
them contains no arithmetic questions, and the interval is honest about
sampling error and silent about the thing that will actually bite you. It is
also at its worst on extremes \dash{} a maximum, a $99$th percentile
\dash{} because the resample can only ever reach values already present.
\end{rigourbox}

One more decision before the technique is usable, and Program~\ref{prog:P25}
named it as a failure mode before this book had said what a bootstrap was.
Your evaluation has $\val{p27.clust.n}$ rows, but they come from
$\val{p27.clust.users}$ users with $\val{p27.clust.rows}$ prompts each. What
do you resample?
\dotline
\nextframe
\end{fr}

\begin{fr}
\ans{The user, not the row}

And getting it wrong is not a rounding error. Program~\ref{prog:P25}'s own
identity is the reason: the variance of a sum of \emph{independent} things
grows like the count, and the variance of a sum of \emph{identical} things
grows like the count squared. Rows from one user are somewhere between the
two, and at the extreme they are one observation wearing
$\val{p27.clust.rows}$ labels.

A bootstrap over rows believes it has $\val{p27.clust.n}$ independent
observations. If the rows within a user agree entirely, it has
$\val{p27.clust.users}$, and the interval it reports is too narrow by a
factor of $\val{p27.clust.factor}$ \dash{} the square root of
$\val{p27.clust.rows}$.

\mermaidfig{p27-what-to-resample}{The unit of resampling is a decision, and
  nothing checks it}{the unit}

Nothing in any library will tell you. The function takes an array, and an
array of $\val{p27.clust.n}$ rows looks exactly like $\val{p27.clust.n}$
observations.
\end{fr}

\begin{fr}
\begin{trapbox}
\textbf{\enquote{Run more resamples and the interval gets tighter.}}

It does not, and the belief is worth killing because resamples are cheap and
running more of them feels like doing something.

There are two different errors in play. The \emph{width} of the interval is
set by how much information the evaluation set contains, which is $n$ items
and nothing else. The number of resamples controls only how precisely you
have located the ends of that interval \dash{} the Monte Carlo error, which
is Program~\ref{prog:P25}'s $1/\sqrt{B}$ in the number of resamples $B$, and
which you can drive as low as you like for the price of arithmetic.

So more resamples buy a more exactly reported interval of \textbf{the same
width}. Buying a narrower one costs items, at
Program~\ref{prog:P25}'s four-times-for-half.
\end{trapbox}

You can now put an interval on anything you can compute. That still leaves
the question the program opened with, and it turns out that asking it
directly is far cheaper than putting an interval on each model separately.
\end{fr}

% ============================================================
\section{They saw the same prompts}
\label{sec:P27-paired}
\index{evaluation!paired comparison}
\index{measurement!sizing a comparison}

\begin{fr}
Program~\ref{prog:P25} answered the sizing question and the answer was
discouraging: to resolve a gap of one point between two models takes
$\val{p27.n.rho00}$ evaluation items \emph{per model}, which is a week of
somebody's budget and more prompts than most benchmarks contain.

That number is right, and it contains an assumption which is almost never
true of the thing you are actually doing. Read the sentence again and name
it.
\dotline
\nextframe
\end{fr}

\begin{fr}
\begin{ansblock}
That the two evaluations are independent of each other.
\end{ansblock}

They are not. You ran both models on \textbf{the same prompts}, because that
is the only sane way to compare them and because collecting a second
evaluation set would cost twice as much.

Program~\ref{prog:P25}'s section 1 has the identity that fixes it, and it
already carried the term this needs:

\[ \Var(X - Y) \;=\; \Var(X) + \Var(Y) - 2\Cov(X, Y) \]

P25 used the independent case, where the covariance is zero and the two
variances add \dash{} which is where its factor of two came from. Here the
covariance is not zero. Say what that does to the item count, and say it as
a formula rather than as a direction.
\dotline
\nextframe
\end{fr}

\begin{fr}
\begin{ansblock}
It multiplies it by $1 - \rho$, where $\rho$ is the correlation between the
two models' per-item results.
\end{ansblock}

Both models sit near the same accuracy $p$, so each per-item indicator has
variance $p(1-p)$ and their covariance is $\rho\,p(1-p)$. The difference then
carries $2p(1-p)(1-\rho)$ where the independent case carried $2p(1-p)$, and
the whole of Program~\ref{prog:P25}'s calculation goes through unchanged with
that one factor attached:

\begin{center}
\begin{tabular}{lrr}
\toprule
$\rho$ & items per model & against $\rho = 0$ \\
\midrule
$0$ & $\val{p27.n.rho00}$ & $1\times$ \\
$\num{0.5}$ & $\val{p27.n.rho50}$ & $\tfrac{1}{2}\times$ \\
$\num{0.8}$ & $\val{p27.n.rho80}$ & $\tfrac{1}{5}\times$ \\
$\num{0.9}$ & $\val{p27.n.rho90}$ & $\tfrac{1}{10}\times$ \\
\bottomrule
\end{tabular}
\end{center}

The script asserts the factor of $1-\rho$ at four accuracies and three gap
sizes rather than the four numbers in the table, because the factor is the
result and the numbers are one cell of it. And the top row is
Program~\ref{prog:P25}'s own committed figure: the paired formula was written
to contain it, and the build fails if the two ever disagree.
\end{fr}

\begin{fr}
So the question is how large $\rho$ is in practice, and the answer is that
for two models worth comparing it is usually large. They are both trained on
most of the internet, they both get the easy items right, and they both fail
the item whose question is ambiguous. Disagreement is the exception, which is
exactly what makes the comparison cheap.

\mermaidfig{p27-same-prompts}{Where the saving in a paired comparison comes
  from}{pairing subtracts}

That is a change of an order of magnitude in what an experiment costs, from
one line of Program~\ref{prog:P25}'s own algebra. It is also the reason the
advice \enquote{evaluate both models on the same prompts} is not merely
tidiness.

Now sharpen it. Of your $\val{p27.items}$ items, some the models both got
right, some they both got wrong, and some exactly one got right. Which of
those three groups carries information about the \emph{difference}?
\dotline
\nextframe
\end{fr}

\begin{fr}
\ans{Only the third}

An item both models answered correctly contributes equally to both scores and
cancels out of the difference exactly. So does an item both failed. Only the
\textbf{discordant} items \dash{} where exactly one model was right \dash{}
say anything at all about which model is better.

Suppose $\val{p27.disc}$ of the $\val{p27.items}$ are discordant, which is
what a high $\rho$ looks like: the two models agreed on
$\val{p27.concord}$ items. Section 1 established that B is ahead by
$\val{p27.net}$ item.

It is worth being explicit about what \enquote{equally good} is being taken
to mean, because the exactness of what follows rests on it. It is not a claim
that the two models are the same model. It is the weaker claim that on an item
where they disagree, either one could have been the one to get it \dash{} so
each discordant item is a coin flip about which of them happened to.

Now the question, and it needs no statistics vocabulary. If the two models
were equally good in that sense, the discordant items would fall to one or the
other like fair coin flips. You flip $\val{p27.disc}$ coins. How surprised
should you be to find them $\val{p27.net}$ away from an even split?
\dotline
\nextframe
\end{fr}

\begin{fr}
\ans{Not at all surprised}

\transcript{p27-thirty-flips}

An even split is the \emph{single most likely} outcome and it still happens
only $\val{p27.p.tie.pct}$ per cent of the time. Everything else is at least
$\val{p27.net}$ away from even, so a result at least this lopsided has
probability $\val{p27.p.exact}$ \dash{} it is what nearly seven flips in
eight of thirty coins look like.

That is Program~\ref{prog:P12}'s binomial coefficient over integers and
nothing else: no normal approximation, no variance, no square root. It is
exact and it is identical on every machine.

So the published difference between the two models is what you would expect
to see if the two models were the same.

Which invites the question in the other direction, and it costs nothing to ask
because it is the same sum run upwards instead of once. Rather than checking
whether your result clears a bar, work out where the bar is.

How large would the gap have to be before it stopped being ordinary?
\dotline
\nextframe
\end{fr}

\begin{fr}
\ans{$\val{p27.net.needed}$ items, which is $\val{p27.net.needed.pts}$ points}

Found by searching upwards from $\val{p27.net}$ rather than quoted: on
$\val{p27.disc}$ discordant items the net has to reach
$\val{p27.net.needed}$ before a fair coin would produce it less than one time
in twenty.

\textbf{On $\val{p27.items}$ items, with these two models, a gap under
$\val{p27.net.needed.pts}$ points is not a result.} That is a number the
reader can carry into a meeting, and it is eleven times the gap that was
published.

Two honest qualifications, both of which make the argument stronger rather
than weaker. The threshold of one in twenty is a convention and the next
section is about what it does and does not mean. And $\val{p27.disc}$
discordant items was assumed rather than measured \dash{} it is a property of
your two models, it is printed by no evaluation harness, and it is the single
most useful number you could add to one.
\end{fr}

\begin{fr}
\begin{aibox}
\textbf{Where this shows up.} Two lines your harness could print and does
not.

Every evaluation report in this field prints two accuracies. Almost none
prints the \textbf{discordant count} and the \textbf{net}, and those two
integers are the entire content of the comparison \dash{} the accuracies are
a lossy summary of them. With them, the arithmetic above is one call to
\code{math.comb} and needs no library at all.

The cost of adding them is a pair of counters in the loop that already
compares the two models item by item. The cost of not adding them is that
the comparison is reported in a form from which it cannot be checked.
\end{aibox}
\end{fr}

% ============================================================
\section{What a p-value says}
\label{sec:P27-p-value}
\index{p-value}
\index{inference!what a p-value does not say}

\begin{fr}
The number $\val{p27.p.exact}$ from the last section has a name, and you have
now computed one before meeting the word, which is the right order.

It is a \textbf{p-value}, and here is the definition, stated flatly enough to
be checked against what you did:

\begin{center}
the probability of a result at least this extreme,\\
\emph{if} the two models were equally good
\end{center}

Every clause is load-bearing. \enquote{At least this extreme} is why the sum
ran over every outcome as lopsided as the one you saw, rather than the one
you saw alone. And the \emph{if} is a condition, which makes the whole thing
a conditional probability \dash{} Program~\ref{prog:P23}'s subject.

Write it as one, using that program's notation, and then say what is on which
side of the bar.
\dotline
\nextframe
\end{fr}

\begin{fr}
\begin{ansblock}
$\Prob(\text{a result this extreme} \mid \text{the models are equally good})$
\dash{} the data is the event, and the claim about the models is the
condition.
\end{ansblock}

Which is the way round nobody reads it.

Suppose a comparison comes back with $p = \num{0.03}$. Is there a three per
cent chance that the two models are equally good?
\dotline
\nextframe
\end{fr}

\begin{fr}
\ans{No, and the two are not related without a third number}

\begin{trapbox}
\textbf{\enquote{$p = \num{0.03}$, so there is a $3$ per cent chance it was a
fluke.}}

That reads $\Prob(A \mid B)$ as $\Prob(B \mid A)$, which is exactly the
inversion Program~\ref{prog:P23} spends a section on, and it is the same
error as the base rate. A p-value conditions on the models being equal; the
sentence conditions on the data. Program~\ref{prog:P23}'s own theorem says
you cannot turn one into the other without a prior, and a p-value comes with
no prior at all.

The reasoning that produces it is not stupid, which is why it is worth
naming rather than scolding. \enquote{This would be unlikely if nothing were
going on, so something is going on} is sound as far as it goes, and it is how
a medical test is read. What it needs, and what a p-value does not supply,
is how often something \emph{is} going on \dash{} the prevalence, in
Program~\ref{prog:P23}'s worked example.
\end{trapbox}
\end{fr}

\begin{fr}
Give it the prior, then, and see what the number becomes.
Program~\ref{prog:P23}'s machinery applies unchanged: put the comparisons
somebody runs in a year on one side, and ask what fraction of the ones that
come back \enquote{significant} are real.

The objection to putting a number there is that it is a guess, and it is
worth answering rather than conceding. It is a guess you are already making
implicitly every time you read a result, and it is one your own team can
estimate: of the changes proposed last year, how many turned out to help? That
is a count somebody could go and make, which puts the prior in the same class
as every other quantity in this program.

Say $\val{p27.prior.real.pct}$ per cent of the comparisons attempted involve a
genuine difference, and say the test finds a genuine one
$\val{p27.power.pct}$ per cent of the time. Of the comparisons that report
$p < \val{p27.alpha}$, what fraction are wrong?
\dotline
\nextframe
\end{fr}

\begin{fr}
\ans{$\val{p27.fdr.pct}$ per cent}

Not $\val{p27.alpha}$ of them. Nearly half.

The arithmetic is two counts and a division, which is
Program~\ref{prog:F10}'s. Out of a hundred comparisons,
$\val{p27.prior.real.pct}$ are real and half of those are found, so five true
findings; ninety are not real and one in twenty of those reports
significance anyway, so four and a half false ones. Four and a half out of
nine and a half.

\textbf{The threshold controls the false positives among the nulls. It says
nothing about the false positives among the findings}, and the second is the
quantity anybody reading a result cares about. Moving the prior moves the
answer and the p-value does not move at all.
\end{fr}

\begin{fr}
One more property, and it is the one that makes the threshold mean anything
at all. If the two models really are identical, and you run the comparison,
how often does $p$ come out below $\val{p27.alpha}$?
\dotline
\nextframe
\end{fr}

\begin{fr}
\ans{$\val{p27.null.rate.pct}$ per cent of the time}

By construction, and that is not a defect \dash{} it is the definition of the
threshold, seen from the other end. \textbf{One comparison in twenty between
two identical models will report a significant difference}, and there is no
version of the procedure in which that is not so.

The measurement here is $\val{p27.null.rate.pct}$ rather than exactly
$\val{p27.alpha}$, and the reason is worth having: the discordant test is
discrete, so its p-value cannot land on the threshold and the achievable
values step over it. A discrete test is therefore slightly
\emph{conservative}, and the script enumerates every one of the
$\val{p27.disc}$ outcomes rather than asserting the continuous answer.

So the p-value is a statement about the procedure's long-run behaviour and
not about your result. It is a useful statement. It is not the one you wanted.
\end{fr}

\begin{fr}
\begin{rigourbox}
\textbf{What you wanted, and where it lives.}

The question in anybody's head is \enquote{what is the probability that B is
better than A}, and no p-value answers it, because that probability requires
a distribution over the claim rather than over the data.

Program~\ref{prog:P28} computes it. Putting a distribution on the parameter,
conditioning on the evaluation, and reading off $\Prob(B > A)$ is the whole
of that program's Beta--Binomial section, and it also contrasts the interval
that comes out with the one this program builds, because the two answer
different questions and are routinely quoted as though they answered the
same one.

This program stops at the frequentist reading, deliberately: it is what the
tools in front of you report, and understanding what it does not say is a
prerequisite for wanting the other thing.
\end{rigourbox}
\end{fr}

% ============================================================
\section{Forty models on one leaderboard}
\label{sec:P27-many}
\index{inference!multiple comparisons}
\index{evaluation!leaderboard}

\begin{fr}
Everything so far has been about one comparison. A leaderboard is not one
comparison.

Put $\val{p27.models}$ models on the same evaluation set and the interesting
quantities are all pairwise, but take the easiest version first: each model
is compared against a baseline, so $\val{p27.models}$ comparisons. Suppose
\textbf{none} of the $\val{p27.models}$ differs from the baseline at all.

What is the chance that at least one of them reports
$p < \val{p27.alpha}$?
\dotline
\nextframe
\end{fr}

\begin{fr}
\ans{$\val{p27.any.false.pct}$ per cent}

$1 - \num{0.95}^{\val{p27.models}}$, which is Program~\ref{prog:F10}'s
complement rule and Program~\ref{prog:P23}'s independence, and nothing else.

\textbf{On a leaderboard of $\val{p27.models}$ models where nothing is
happening, something looks significant almost every time.} The threshold was
built to be crossed one time in twenty by accident, and forty attempts is
forty chances to be that one.

The crude repair is to divide the threshold by the number of comparisons
\dash{} $\val{p27.alpha}$ becomes $\val{p27.alpha}/\val{p27.models}$ \dash{}
which is due to Bonferroni and is what most fields do. It works, in the sense
that it restores the overall rate. What does it cost?
\dotline
\nextframe
\end{fr}

\begin{fr}
\ans{$\val{p27.bonf.cost}$ times the items}

The threshold moves the multiplier in Program~\ref{prog:P25}'s interval from
$\val{p27.z.bonf}$ in place of $\num{1.96}$, and section 3's formula carries
that multiplier \emph{squared}. So a leaderboard that wants forty honest
comparisons needs $\val{p27.bonf.cost}$ times the evaluation items of one,
on top of everything else.

Which is the trade in the open, and it is the right way to see it: multiple
comparisons is not a statistical technicality but a budget line. Ask forty
questions of one evaluation set and either you accept that one of the answers
is wrong, or you buy nearly three times as many items.
\end{fr}

\begin{fr}
There is a second effect of running many comparisons and it is different in
kind, larger, and almost never mentioned. It survives even when nobody
computes a p-value at all.

Take $\val{p27.models}$ models whose true accuracies are \textbf{exactly
equal}, at $\val{p27.lb.p.pct}$ per cent. Evaluate each on
$\val{p27.lb.n}$ items and sort the table. The model
at the top is the one whose evaluation went best, and its observed score sits
above its own true accuracy.

By how much, roughly, in points? The standard error of one score here is
$\val{p27.lb.se}$ points.
\dotline
\nextframe
\end{fr}

\begin{fr}
\ans{About $\val{p27.lb.margin}$ points}

The winner of $\val{p27.models}$ equal contestants sits
$\val{p27.emax}$ standard errors above the truth, on average, and
$\val{p27.emax} \times \val{p27.lb.se}$ is $\val{p27.lb.margin}$.

That figure is the expected maximum of $\val{p27.models}$ standard normal
draws, and it is \textbf{integrated rather than simulated} \dash{} for
Program~\ref{prog:P24}'s reason, that \enquote{it agreed within the error
bar} is exactly the reading a section like this one exists to refuse. The
two cases with closed forms, $m = 1$ and $m = 2$, are checked against them
first.

\mermaidfig{p27-winner-is-noise}{What sorting a table of equals does to the
  top of it}{winner's margin}

So on a leaderboard where a one-point gap is reported as a result, the model
at the top of a field of forty appears about three points better than it is,
\textbf{with no differences between the models at all}. Selection did that,
not measurement.
\end{fr}

\begin{fr}
\begin{trapbox}
\textbf{\enquote{We tried thirty configurations and picked the one that
scored best on the evaluation set.}}

Which is the same arithmetic, said about your own work rather than about a
public table \dash{} and it is the reason a held-out set exists, stated as a
number rather than as a rule somebody handed you.

The chosen configuration's score is not an estimate of how it will do. It is
the maximum of thirty noisy estimates, and the maximum of $m$ things is
biased upwards by an amount that grows with $m$. Reporting it as the expected
performance is reporting the selection along with the model.

The repair is not statistical. Choose on one set and \emph{report} on
another, and the number you report has been selected on nothing.
\end{trapbox}
\end{fr}

% ============================================================
\section{What you can now check}
\label{sec:P27-check}
\index{evaluation!what to ask of a published number}
\index{inference!what this program does not claim}

\begin{fr}
Back to the sentence the program opened with, and it can now be answered
rather than doubted.

\begin{center}
\emph{Model B scored $\val{p27.quoted.hi}$ and model A scored
$\val{p27.quoted.lo}$ on $\val{p27.items}$ evaluation items.}
\end{center}

Four things, and none of them needed a library.

\textbf{The score is not on the grid.} $\val{p27.quoted.hi}$ per cent of
$\val{p27.items}$ items is $\val{p27.k.hi}$ items, so either the denominator
is not $\val{p27.items}$ or the metric is not an accuracy. Ask which.

\textbf{The gap is $\val{p27.gap.items}$ item}, and a percentage is an
excellent way to hide that.

\textbf{The interval on either score is over
$\val{p27.hw.ratio.floor}$ times the gap}, from one square root.

\textbf{And the paired test, which is the right one, gives
$\val{p27.p.exact}$} \dash{} on $\val{p27.disc}$ discordant items, which is
the one number you would have to ask for.

The gap would have to be $\val{p27.net.needed.pts}$ points before this
evaluation set could distinguish the two models at all.
\end{fr}

\begin{fr}
Two things this program does not claim, and both are worth more than another
technique.

\textbf{It has not measured how many published results are noise.} It has
computed what it would take for one of them not to be, on one worked
evaluation, and that is a different statement. Doing the first properly means
collecting the discordant counts behind a large number of published
comparisons, and those are not published \dash{} which is itself the
finding, and it is why section 3's aibox asks for two integers rather than
for better statistics.

\textbf{And none of this is a decision procedure.} A p-value below a
threshold is not permission to ship and one above it is not proof of
equivalence: section 4's own arithmetic says the second especially, because a
test that would need $\val{p27.net.needed.pts}$ points to see anything has
said nothing whatever about a real difference of one.

What is left when a comparison cannot resolve the thing you care about is a
design question rather than an arithmetic one. More items, at
Program~\ref{prog:P25}'s price. The same prompts, at section 3's discount.
Or a different measurement.

Those three are not equally priced and the order is worth having. Pairing is
free and is usually already available. Items cost linearly in money and
quadratically in precision. And a different measurement \dash{} one where the
thing you care about produces a larger signal \dash{} is the only one of the
three with no ceiling, and the only one nobody budgets for.

There is a fourth thing available, and it is the cheapest of all.
\dotline
\nextframe
\end{fr}

\begin{fr}
\begin{ansblock}
Or the honest answer, which is available for nothing: report that the two
models could not be distinguished on this evaluation, and say how large a
difference would have had to be.
\end{ansblock}

That last clause is what turns \enquote{no significant difference} from a
non-result into a measurement, and almost nobody writes it. It is the
$\val{p27.net.needed.pts}$ points from section 3, and computing it costs a
loop over \code{math.comb}.

Where this leaves you. You can read a benchmark number against its own
denominator; you can put an interval on a metric that has no formula, and
name the unit to resample; you can size and run the comparison people
actually want, using the fact that both models saw the same prompts; you can
say what a p-value does and does not say, without either revering it or
dismissing it; and you can price a leaderboard's worth of comparisons before
reading its top row.

Program~\ref{prog:P28} takes the question this program had to hand back
\dash{} the probability that B is better \dash{} and answers it, by putting a
distribution on the thing you are uncertain about rather than on the data you
already have.
\end{fr}

\begin{summarybox}
\sumitem{1--3}{\textbf{Read a score against its own denominator first.} A
  score of $k$ right out of $n$ lands on a grid $100/n$ points apart, so
  $\val{p27.quoted.hi}$ per cent on $\val{p27.items}$ items is
  $\val{p27.k.hi}$ items and is not achievable. A one-decimal accuracy needs
  over $\val{p27.decimal.floor}$ items before the decimal can round.}
\sumitem{3--5}{\textbf{The published gap is $\val{p27.gap.items}$ item.}
  $\val{p27.gap.pts}$ points of $\val{p27.items}$ items, against an interval
  of $\pm\val{p27.hw.200}$ points \dash{} over
  $\val{p27.hw.ratio.floor}$ times the gap.}
\sumitem{6}{Comparing two separate intervals by whether they overlap is not
  a test. Build the interval on the difference, which is the quantity the
  question is about.}
\sumitem{7--9}{\textbf{The bootstrap} resamples the sample with replacement
  and needs no formula for any metric. On a proportion its distribution is
  \emph{exactly} binomial, so it agrees with Program~\ref{prog:P25}'s formula
  to within one item by construction rather than by luck.}
\sumitem{10--11}{Its use is the case with no formula \dash{} a median, whose
  interval $\intcc{\val{p27.med.lo}}{\val{p27.med.hi}}$ has endpoints that
  are measurements somebody took. It cannot invent items you did not
  evaluate, and it is worst on extremes.}
\sumitem{12--14}{\textbf{The unit of resampling is a decision and nothing
  checks it.} $\val{p27.clust.n}$ rows from $\val{p27.clust.users}$ users are
  not $\val{p27.clust.n}$ observations, and a bootstrap over rows reports an
  interval too narrow by $\val{p27.clust.factor}$. More resamples never
  narrow an interval; only items do.}
\sumitem{15--17}{\textbf{Pairing multiplies the items you need by
  $1-\rho$}, exactly, at every accuracy and gap. Program~\ref{prog:P25}'s
  $\val{p27.n.rho00}$ becomes $\val{p27.n.rho90}$ at $\rho = \num{0.9}$
  \dash{} an order of magnitude, from the covariance term in
  $\Var(X-Y)$.}
\sumitem{18--19}{\textbf{Only the discordant items carry the difference.}
  The $\val{p27.concord}$ items both models answered the same way cancel
  exactly, and the comparison is entirely about the
  $\val{p27.disc}$ that are left.}
\sumitem{19--20}{Under \enquote{equally good} those items are fair coin
  flips, so the test is Program~\ref{prog:P12}'s binomial coefficient over
  integers \dash{} exact, machine-independent, and needing no variance. The
  worked gap gives $\val{p27.p.exact}$.}
\sumitem{21}{\textbf{On this evaluation a gap under
  $\val{p27.net.needed.pts}$ points is not a result}, found by search. The
  published gap was $\val{p27.gap.pts}$.}
\sumitem{23--25}{\textbf{A p-value is
  $\Prob(\text{data this extreme} \mid \text{no difference})$}, and reading
  it backwards is Program~\ref{prog:P23}'s inversion \dash{} the same error
  as the base rate.}
\sumitem{26--27}{Given a prior it can be turned round, and the answer is
  alarming: at $\val{p27.prior.real.pct}$ per cent of comparisons real and
  $\val{p27.power.pct}$ per cent power, $\val{p27.fdr.pct}$ per cent of
  \enquote{significant} findings are false.}
\sumitem{28--29}{\textbf{One comparison in twenty between identical models
  reports a difference}, by construction. Here $\val{p27.null.rate.pct}$ per
  cent, because a discrete test steps over the threshold rather than landing
  on it.}
\sumitem{31--33}{\textbf{$\val{p27.models}$ comparisons make that
  $\val{p27.any.false.pct}$ per cent.} Dividing the threshold restores it and
  costs $\val{p27.bonf.cost}$ times the items \dash{} multiple comparisons
  is a budget line, not a technicality.}
\sumitem{34--35}{\textbf{And sorting a table biases its top row.} The winner
  of $\val{p27.models}$ \emph{equal} models sits $\val{p27.emax}$ standard
  errors above its own truth, which is $\val{p27.lb.margin}$ points at
  $\val{p27.lb.n}$ items. Selection, not measurement \dash{} and the same
  arithmetic is why a held-out set exists.}
\sumitem{37--39}{What is left when a comparison cannot resolve what you care
  about is a design question. \textbf{And \enquote{no significant
  difference} becomes a measurement the moment you say how large a difference
  would have had to be.}}
\end{summarybox}

\canyou

\begin{testexercises}
\item A model is reported at $\num{83.3}$ per cent on $60$ evaluation items.
  Is that achievable, and what are the two nearest scores that are?
  \answerto{Yes: $\num{83.3}$ per cent of $60$ is exactly $50$ items. The
  grid here is $100/60 = \num{1.67}$ points, so the neighbours are $49/60 =
  \num{81.7}$ and $51/60 = \num{85.0}$ per cent \dash{} which is how coarse
  a $60$-item benchmark is.}
\item Two models are evaluated on $500$ shared prompts. They disagree on
  $80$ of them, and of those, $46$ go to B and $34$ to A. What is the net, in
  items and in points?
  \answerto{The net is $46 - 34 = 12$ items, which is
  $12/500 = \num{2.4}$ points. Note that the two accuracies differ by
  $\num{2.4}$ points however many items they both got right \dash{} the
  concordant ones cancel.}
\item An evaluation has $600$ rows drawn from $50$ conversations. A colleague
  bootstraps over rows and reports an interval. In which direction is it
  wrong, and by roughly what factor if the rows within a conversation agree
  entirely?
  \answerto{Too narrow, by up to $\sqrt{600/50} = \sqrt{12} \approx
  \num{3.5}$. The bootstrap believes it has $600$ independent observations
  and in the worst case it has $50$.}
\item A team runs $20$ comparisons against a baseline and reports the two
  that came back at $p < \num{0.05}$. If none of the $20$ systems differs
  from the baseline, how many such reports would you expect?
  \answerto{$20 \times \num{0.05} = 1$. Finding two is not evidence of
  anything; the expected number under no difference at all is one, and
  $\Prob(\text{at least one}) = 1 - \num{0.95}^{20} = \num{0.64}$.}
\item Your comparison needs $\val{p27.n.rho00}$ items per model if the two
  evaluations are independent. You have $\val{p27.n.rho80}$ prompts and can
  run both models on all of them. What correlation between the two models'
  per-item results would you need for that to be enough?
  \answerto{$\val{p27.n.rho00} \times (1 - \rho) \le \val{p27.n.rho80}$
  gives $\rho \ge \num{0.8}$ \dash{} which is what the table in section 3
  says, read backwards.}
\end{testexercises}

\begin{furtherproblems}
\item A benchmark reports $\num{62.48}$ per cent. What is the smallest number
  of items on which that many decimal places could be computed at all?
  \answerto{The grid must be finer than half a hundredth of a point, so
  $1/n < \num{0.00005}$ and $n > 20\,000$. Almost no benchmark has that many
  items, so two decimal places on a public leaderboard are decoration.}
\item Show that the bootstrap distribution of a \emph{mean} of $n$ values has
  variance $(1 - 1/n)$ times the usual $s^{2}/n$, and say why that does not
  matter at $n = \val{p27.items}$.
  \answerto{Resampling draws from the empirical distribution, whose variance
  is the $n$-denominator one \dash{} Program~\ref{prog:P26}'s $\frac{n-1}{n}$
  factor. At $n = \val{p27.items}$ that is a $\num{0.5}$ per cent difference
  in the variance and a quarter of a per cent in the interval, which is far
  inside the grid section 1 measured.}
\item Two models are compared on $n$ shared items and disagree on a fraction
  $d$ of them. Write the number of items needed to resolve a gap of $\delta$
  in terms of $d$ rather than $\rho$, and say which of the two an evaluation
  harness could actually report.
  \answerto{The discordant fraction is what the harness can count, and it is
  the quantity the exact test uses directly: the net has to reach about
  $\num{1.96}\sqrt{dn}$ items for a fair coin not to produce it. $\rho$ has
  to be estimated; $d$ is two counters in a loop, which is why section 3's
  aibox asks for it.}
\item A leaderboard's top two models are separated by $\num{0.3}$ points on
  $\val{p27.lb.n}$ items. Without any further information, what is the most
  useful thing you can say?
  \answerto{That the separation is $3$ items, that the standard error on
  either score is $\val{p27.lb.se}$ points, and that with
  $\val{p27.models}$ models in the table the top row is expected to sit
  about $\val{p27.lb.margin}$ points above its own true accuracy from
  selection alone \dash{} which is over ten times the gap being reported.}
\item Program~\ref{prog:P26} showed that a deliberately biased estimator can
  beat an unbiased one. Does the same trade apply to a confidence interval,
  and what would the two axes be?
  \answerto{Yes, and the axes are \emph{coverage} and \emph{width}. An
  interval that is always $\intcc{0}{100}$ has perfect coverage and no use;
  one that is always a single point has zero width and no coverage. The
  conventional interval fixes coverage at $95$ per cent and reports whatever
  width that costs, which is one choice on that trade rather than the only
  one.}
\end{furtherproblems}
